\chapter{The Hartree-Fock Approximation}\label{c3}
\section{An overview}\label{c3s1}
The Hartree-Fock method is an algorithm to estimate the best variational 
solution to the hamiltonian $H$ of a system of $N$-electrons, under the 
Born-Oppenheimer approximaton. The ground state of such a system is described 
by the Slater determinant
\begin{equation}\label{c3e1}
\ket{\Psi_0} = \ket{\chi_1\ldots\chi_a\ldots\chi_N}.
\end{equation}
The algorithm minimises $\braket{\Psi_0|H|\Psi_0}$ subject to the constraint
\begin{equation}\label{c3e2}
\braket{\chi_a|\chi_b} = \delta_{ab}.
\end{equation}
From equation \eqref{c2e102},
\begin{equation}\label{c3e3}
\braket{\Psi_0|H|\Psi_0} = \sum_{a=1}^N[a|h|a] + \frac{1}{2}\sum_{a=1}^N
\sum_{b=1}^N\left([aa|bb] - [ab|ba]\right).
\end{equation}
Thus, we minimise \eqref{c3e3} subject to the constraint \eqref{c3e2}. The
spin orbitals of the minimal energy state satisfy the equation
\begin{equation}\label{c3e4}
\begin{aligned}
h(\vec{x}_1)\chi_a(\vec{x}_1) + 
\sum_{b\ne a}\left[\int d\vec{x}_2\chi_b^\ast(\vec{x}_2)\chi_b(\vec{x}_2)
r_{12}^{-1}\right]\chi_a(\vec{x}_1)& \\
 - \sum_{b\ne a}\left[\int d\vec{x}_2\chi_b^\ast(\vec{x}_2)\chi_a(\vec{x}_2)
r_{12}^{-1}\right]\chi_b(\vec{x}_1) 
& = \varepsilon_a\chi_a(\vec{x}_1),
\end{aligned}
\end{equation}
where
\begin{equation}\label{c3e5}
h(\vec{x}_1) = -\frac{1}{2}\nabla_1^2 - \sum_A\frac{Z_A}{r_{1A}},
\end{equation}
is the hamiltonian of electron 1 in the field of all nuclei. The term
\[
\int d\vec{x}_2\chi_b^\ast(\vec{x}_2)\chi_b(\vec{x}_2)r_{12}^{-1} =
\int d\vec{x}_2\frac{\rho_b(\vec{x}_2)}{r_{12}}
\]
is the Coulomb potential due to the $b$th electron at $\vec{x}_1$. When it is
summed over all $b \ne a$, it is the electrostatic potential due to the $N-1$
electrons at the position $\vec{x}_1 = (\vec{r}_1, \omega_1)$ of electron 1.
(Remember that in atomic units, $e = 1$.) Therefore,
\[
\int d\vec{x}_2\chi_b^\ast(\vec{x}_2)\chi_b(\vec{x}_2)r_{12}^{-1}
\chi_a(\vec{x}_1)
\]
is the electrostatic, repulsive potential energy of electron 1 in the field of
the remaining $N-1$ electrons. The third term in \eqref{c2e4} is a bit tricky.
We introduce an operator $K_b$ as
\begin{equation}\label{c3e6}
K_b\chi_a(\vec{x}_1) = \int d\vec{x}_2 \chi_b^\ast(\vec{x}_2)\chi_a(\vec{x}_2)
r_{12}^{-1}\chi_b(\vec{x}_1).
\end{equation}
This equation is very similar to the Coulomb term except that the electrons
1 and 2 are swapped. The Coulomb term had electron 1 in spin orbital $\vec{x}_1$
and electron 2 in $\vec{x}_2$. In equation \eqref{c2e6} it is the other way 
round. If we introduce the electron swapping operator
\begin{equation}\label{c3e7}
P_{12}\chi_a(\vec{x}_1)\chi_b(\vec{x}_2) = \chi_a(\vec{x}_2)\chi_b(\vec{x}_2)
\end{equation}
then we can write \eqref{c3e6} as
\begin{equation}\label{c3e8}
\begin{aligned}
K_b\chi_a(\vec{x}_1) &= \int d\vec{x}_2 \chi_b^\ast(\vec{x}_2)
r_{12}^{-1}P_{12}\chi_a(\vec{x}_1)\chi_b(\vec{x}_2) \\
 &= \int d\vec{x}_2\chi_b^\ast(\vec{x}_2) r_{12}^{-1}P_{12}\chi_b(\vec{x}_2)
 \chi_a(\vec{x}_1).
\end{aligned}
\end{equation}
The operator $K_b$ is itself a function of $\vec{x}_1$. We can define it as
\begin{equation}\label{c3e9}
K_b(\vec{x}_1) = 
\int d\vec{x}_2\chi_b^\ast(\vec{x}_2)r_{12}^{-1}P_{12}\chi_b(\vec{x}_2).
\end{equation}
It is called the \emph{exchange operator}. The analogous operator
\begin{equation}\label{c3e10}
J_b(\vec{x}_1) = 
\int d\vec{x}_2\chi_b^\ast(\vec{x}_2)r_{12}^{-1}\chi_b(\vec{x}_2)
\end{equation}
is called the \emph{Coulomb operator}. The terms $K_b(\vec{x}_1)\chi_a(\vec{x}_1)$
and $J_b(\vec{x}_1)\chi_a(\vec{x}_1)$ are called the \emph{exchange energy} and
\emph{Coulomb energy}. In terms of these, equation \eqref{c3e4} is
\begin{equation}\label{c3e11}
h(\vec{x}_1)\chi_a(\vec{x}_1) + \sum_{b \ne a}\left[J_b(\vec{x}_1) - 
K_b(\vec{x}_1)\right]\chi_a(\vec{x}_1) = \varepsilon_a\chi_a(\vec{x}_1).
\end{equation}
We can simplify it further by defining the Fock operator
\begin{equation}\label{c3e12}
f(1) = h(1) + \sum_{b \ne a}(J_b(1) - K_b(1))
\end{equation}
so that we have
\begin{equation}\label{c3e13}
f(\vec{x}_1)\chi_a(\vec{x}_1) = \varepsilon_a\chi_a(\vec{x}_1).
\end{equation}
In spite of its appearance, \eqref{c3e13} is not a linear eigenvalue equation.
This is because the operator $f(\vec{x}_1)$ depends on $\chi_a(\vec{x}_1)$
through the Coulomb and exchange terms.

From \eqref{c3e6} it is evident that the term $K_b(\vec{x}_1)\chi_a(\vec{x}_1)$
depends on values of $\chi_a$ at all points, not just $\vec{x}_1$. It is,
therefore, a non-local operator. 

The expectation values of these operators are
\begin{align}
\braket{\chi_a|J_b|\chi_a} &= \iint d\vec{x}_1d\vec{x}_2\chi_a^\ast(\vec{x}_1)
\chi_a(\vec{x}_1)r_{12}^{-1}\chi_b^\ast(\vec{x}_2)\chi_b(\vec{x}_2) = [aa|bb] 
\label{c3e14} \\
\braket{\chi_a|K_b|\chi_a} &= \iint d\vec{x}_1d\vec{x}_2\chi_a^\ast(\vec{x}_1)
\chi_b(\vec{x}_1)r_{12}^{-1}\chi_b^\ast(\vec{x}_2)\chi_a(\vec{x}_2) = [ab|ba] 
\label{c3e15}
\end{align}

\section{Derivation of Hartree-Fock equation}\label{c3s2}
We will derive equation \eqref{c3e4} by minimising \eqref{c3e3} subject to the
constraint that the spin-orbitals are orthonormal. Consider the determinant
\[
\ket{\Psi_0} = \ket{\chi_1\ldots\chi_a\ldots\chi_N}
\]
such that
\begin{equation}\label{c3e16}
\braket{\chi_a|\chi_b} = [a|b| = \delta_{ab} \Rightarrow [a|b] - \delta_{ab}=0.
\end{equation}
Minimisation of $\braket{\Psi_0|H|\Psi_0}$ subject to the constraint 
\eqref{c3e16} involves finding the extrema of
\begin{equation}\label{c3e17}
\mathscr{L}[\{\chi_a\}] = \braket{\Psi_0|H|\Psi_0} - \sum_{ab}\varepsilon_{ab}(
[a|b] - \delta_{ab}),
\end{equation}
where the $N^2$ constants $\varepsilon_{ab}$ are Lagrange's undetermined 
multipliers. From equation \eqref{c2e102},
\begin{equation}\label{c3e18}
\mathscr{L}[\{\chi_a\}] = \sum_{a=1}^N[a|h|a] + \frac{1}{2}\sum_{a=1}^N
\sum_{b=1}^N([aa|bb] - [ab|ba]) - \sum_{ab}\varepsilon_{ab}([a|b]-\delta_{ab}).
\end{equation}
The first variation of $\mathscr{L}$ is
\begin{equation}\label{c3e19}
\begin{aligned}
\delta\mathscr{L} &= \sum_{a=1}^N[\delta a|h|a] + \sum_{a=1}^N[a|h|\delta a] +\\
 &\quad\frac{1}{2}\sum_{a,b=1}^N\left([\delta aa|bb] + [a\delta a|bb] + 
 [aa|\delta bb] + [aa|b\delta b]\right) - \\
 &\quad\frac{1}{2}\sum_{a,b=1}^N\left([\delta ab|ba] + [a\delta b|ba] + 
 [ab|\delta ba] + [ba|b\delta a]\right) - \\
 &\quad\sum_{a,b=1}^N\varepsilon_{ab}([\delta a|b] + [a|\delta b]).
\end{aligned}
\end{equation}
Now,
\[
[a|h|b] = \braket{a|h|b} = \int d\vec{x}\chi_a^\ast(\vec{x})h(\vec{r})\chi_b(
\vec{x})
\]
so that
\begin{equation}\label{c3e20}
[\delta a|h|a] = [a|h|\delta a]^\ast,
\end{equation}
because $h$ is real. Therefore,
\begin{equation}\label{c3e21}
\sum_{a=1}^N[\delta a|h|a] + \sum_{a=1}^N[a|h|\delta a] =
\sum_{a=1}^N[\delta a|h|a] + \text{c.c.},
\end{equation}
where c.c. stands for complex conjugate.
Since
\[
[ad|cd] = \iint d\vec{x}_1d\vec{x}_2 \chi_a^\ast(\vec{x}_1)\chi_b(\vec{x}_1)
r_{12}^{-1}\chi_c^\ast(\vec{x}_2)\chi_d(\vec{x}_2),
\]
we have $[(\delta a)a|bb] = [a(\delta a)|bb]^\ast, [aa|(\delta b)b] = 
[aa|b(\delta b)]^\ast$. We can always interchange the variables of integration, 
so that $[(\delta a)a|bb] = [bb|(\delta a)a] = [aa|(\delta b)b]$. Therefore
\begin{equation}\label{c3e22}
\frac{1}{2}\sum_{a,b=1}^N\left([(\delta a)a|bb] + [a(\delta a)|bb] + 
 [aa|(\delta b)b] + [aa|b(\delta b)]\right) = \sum_{a,b=1}^N[(\delta a)a|bb] +
 \text{c.c.}
\end{equation}
Similarly
\begin{equation}\label{c3e23}
\frac{1}{2}\sum_{a,b=1}^N\left([(\delta a)b|ba] + [a(\delta b)|ba] + 
 [ab|(\delta b)a] + [ba|b(\delta a)]\right) =
\sum_{a,b=1}^n[(\delta a)b|ba] + \text{c.c.}
\end{equation}
Lastly,
\[
\begin{aligned}
\sum_{a,b=1}^N\varepsilon_{ab}([\delta a|b] + [a|\delta b]) &=
\sum_{a,b=1}^N\varepsilon_{ab}[\delta a|b] + \sum_{a,b=1}^N\varepsilon_{ab}
[\delta b|a]^\ast \\
 &= \sum_{a,b=1}^N\varepsilon_{ab}[\delta a|b] + \sum_{a,b=1}^N\varepsilon_{ba}
[\delta a|b]^\ast
\end{aligned}
\]
Since $\mathscr{L}$ is real, $\varepsilon_{ab} = \varepsilon_{ba}^\ast$ and 
hence
\begin{equation}\label{c3e24}
\sum_{a,b=1}^N\varepsilon_{ab}([\delta a|b] + [a|\delta b]) =
\sum_{a,b=1}^N\varepsilon_{ab}[\delta a|b] + \text{c.c.}
\end{equation}
Substituting \eqref{c3e21} to \eqref{c3e24} into \eqref{c3e19}, we get
\begin{equation}\label{c3e25}
\delta\mathscr{L} = \sum_{a=1}^N[\delta a|h|a] + \sum_{a,b=1}^N\left([(\delta a)
a|bb] - [(\delta a)b|ba]\right) - \sum_{a,b=1}^N\varepsilon_{ab}[(\delta a)|b]
+ \text{c.c.}
\end{equation}
Writing the terms as integrals and setting the first variation to zero gives,
\[
\begin{aligned}
\delta\mathscr{L} &=
\sum_{a=1}^N\int d\vec{x}_1\delta\chi_a^\ast(\vec{x}_1)\left[h(\vec{x}_1)
\chi_a(\vec{x}_1) + \sum_{b=1}^N\int d\vec{x}_2\chi_a(1)r_{12}^{-1}
\chi_b^\ast(2)\chi_b(2) - \right. \\
&\quad\left.\sum_{b=1}^N\int d\vec{x}_2 \chi_b(1)r_{12}^{-1} \chi_b^\ast(2)\chi_a(2)
- \sum_{b=1}^N \varepsilon_{ab}\chi_b(1) \right] + \text{c.c.}\\
&= \sum_{a=1}^N\int d\vec{x}_1\delta\chi_a^\ast(\vec{x}_1)\left[h(\vec{x}_1)
\chi_a(\vec{x}_1) + \sum_{b=1}^N\int d\vec{x}_2\chi_a(1)r_{12}^{-1}
\chi_b^\ast(2)\chi_b(2) - \right. \\
&\quad\left.\sum_{b=1}^N\int d\vec{x}_2 \chi_b^\ast(2)r_{12}^{-1} \chi_b(1)\chi_a(2)
- \sum_{b=1}^N\varepsilon_{ab}\chi_b(1) \right]+ \text{c.c.} \\
&= \sum_{a=1}^N\int d\vec{x}_1\delta\chi_a^\ast(\vec{x}_1)\left[h(\vec{x}_1)
\chi_a(\vec{x}_1) + \sum_{b=1}^N\int d\vec{x}_2\chi_b^\ast(2)r_{12}^{-1}
\chi_b(2)\chi_a(1) - \right. \\
&\quad\left.\sum_{b=1}^N\int d\vec{x}_2 \chi_b^\ast(2)r_{12}^{-1}P_{12} 
\chi_b(2)\chi_a(1) - 
\sum_{b=1}^N\varepsilon_{ab}\chi_b(1) \right] + \text{c.c.}\\
&= \sum_{a=1}^N\int d\vec{x}_1\delta\chi_a^\ast(1)\left[\left(h(1) + \sum_{b=1}^N(
J_b(1) - K_b(1))\right)\chi_a(1) - \right. \\
&\quad\left.\sum_{b=1}^N\varepsilon_{ab}\chi_b(1) \right] 
+ \text{c.c.}
\end{aligned}
\]
The variations $\delta\chi_a^\ast$ (and $\delta\chi_a$) are arbitrary so that
$\delta\mathscr{L} = 0$ implies
\begin{equation}\label{c3e26}
\left(h(1) + \sum_{b=1}^N( J_b(1) - K_b(1))\right)\chi_a(1) = \sum_{b=1}^N
\varepsilon_{ab}\chi_b(1)
\end{equation}
and its complex conjugate. The term
\begin{equation}\label{c3e27}
v^{HF} = \sum_{b=1}^N( J_b(1) - K_b(1))
\end{equation}
is the effective potential felt by an electron due to the other $N-1$ electrons.
It is called the \emph{Hartree-Fock} potential.

Let $U$ be a unitary transformation that takes the spin orbitals $\{\chi_a\}$
to $\{\chi_a^\op\}$. Thus,
\begin{equation}\label{c3e28}
\chi_a^\op = \sum_{b}U_{ab}\chi_b
\end{equation}
and
\begin{equation}\label{c3e29}
\chi_a^{\op\ast} = \sum_b U_{ab}^\ast \chi_b^\ast.
\end{equation}
Since $U$ is unitary, $UU^\dagger = U^\dagger U =  I$, that is
\begin{equation}\label{c3e30}
\begin{aligned}
\sum_b U_{ab}U^\dagger_{bc} &= \sum_b U_{ab}U_{cb}^\ast = \delta_{ac} \\
\sum_b U^\dagger_{ab}U_{bc} &= \sum_b U_{ba}^\ast U_{bc} = \delta_{ac}.
\end{aligned}
\end{equation}
If $J_b^\op(1)$ and $K_b^\op(1)$ are the Coulomb and the
exchange operators in the primed basis then,
\begin{equation}\label{c3e31}
\begin{aligned}
\sum_b J_b^\op(1) &= 
\sum_b \int d\vec{x}_2\chi_b^{\op\ast}(2)r_{12}^{-1}\chi_b^\op(2) \\
 &=\sum_b 
 \int d\vec{x}_2\sum_{c,d}U_{bc}^\ast\chi_c^\ast(2)r_{12}^{-1}U_{bd}\chi_d(2) \\
 &=\sum_{c,d} 
 \int d\vec{x}_2\sum_bU_{bc}^\ast U_{bd}\chi_c^\ast(2)r_{12}^{-1}\chi_d(2) \\
 &= \sum_{c,d} \int d\vec{x}_2\delta_{cd}\chi_c^\ast(2)r_{12}^{-1}\chi_d(2) \\
 &= \sum_c \int d\vec{x}_2 \chi_c^\ast(2)r_{12}^{-1}\chi_c(2) \\
 &= \sum_c J_c(1) = \sum_b J_b(1)
\end{aligned}
\end{equation}
and
\begin{equation}\label{c3e32}
\begin{aligned}
\sum_b K_b^\op(1) &= 
\sum_b \int d\vec{x}_2\chi_b^{\op\ast}(2)r_{12}^{-1}P_{12}\chi_b^\op(2) \\
 &=\sum_b \int d\vec{x}_2\sum_{c,d}U_{cb}^\ast\chi_c^\ast(2)r_{12}^{-1}
   P_{12}U_{bd}\chi_d(2) \\
 &=\sum_{c,d} \int d\vec{x}_2\sum_bU_{cb}^\ast U_{bd}\chi_c^\ast(2)r_{12}^{-1}
   P_{12}\chi_d(2) \\
 &= \sum_{c,d} 
    \int d\vec{x}_2\delta_{cd}\chi_c^\ast(2)r_{12}^{-1}P_{12}\chi_d(2) \\
 &= \sum_c \int d\vec{x}_2 \chi_c^\ast(2)r_{12}^{-1}P_{12}\chi_c(2) \\
 &= \sum_c K_c(1) = \sum_b K_b(1)
\end{aligned}
\end{equation}
Thus, the Coulomb and exchange operators are invariant under a unitary 
transformation. The term $h(1)$ is independent of spin orbitals and is unaffected
by a unitary transformation of the spin orbitals. Therefore, if
\begin{equation}\label{c3e33}
f(1) = h(1) + \sum_{b=1}^N( J_b(1) - K_b(1)
\end{equation}
is the \emph{Fock operator} then $f(1) = f^\op(1)$. In terms of the Fock 
operator, equation \eqref{c3e26} is
\begin{equation}\label{c3e34}
f(1)\chi_a(1) = \sum_{b=1}^N \varepsilon_{ab}\chi_b(1).
\end{equation}
so that
\[
\braket{\chi_c|f|\chi_a} = \sum_{b=1}^N\varepsilon_{ab}\braket{\chi_c|\chi_b}.
\]
Using the orthornormality constraint of \eqref{c3e19}, we get
\begin{equation}\label{c3e35}
\varepsilon_{ca} = \braket{\chi_c|f|\chi_a}.
\end{equation}
Under a unitary transformation it becomes,
\begin{equation}\label{c3e36}
\begin{aligned}
\varepsilon_{ca}^\op &= \braket{\chi_c^\op|f^\op|\chi_a^\op} \\
 &= \int d\vec{x}_1 \chi_c^{\op\ast}(1) f(1) \chi_a^\op(1) \\
 &= \sum_{d,e} \int d\vec{x}_1 U^\ast_{cd}\chi_d^\ast(1)f(1) U_{ae}\chi_e(1) \\
 &= \sum_{d, e} U^\ast_{cd} \int d\vec{x}_1\chi_d^\ast(1)f(1)\chi_e(1)U_{ae} \\
 &= \sum_{d, e} U^\ast_{cd} \varepsilon_{de}U_{ae} \\
 &= \sum_{d, e} U^\dagger_{dc} \varepsilon_{de} U_{ae}
\end{aligned}
\end{equation}
Thus, the Lagrange multipliers vary as
\begin{equation}\label{c3e37}
\varepsilon^\op = U^\dagger\varepsilon U.
\end{equation}

Equation \eqref{c3e34} can be written in matrix form as 
\[
f(1)\ket{\chi_a} = \sum_{b}\varepsilon_{ab}\ket{\chi_b},
\]
so that under a unitary transformation that diagonalises the matrix 
$\varepsilon$, it becomes
\begin{equation}\label{c3e38}
f^\op(1)\ket{\chi_a^\op} = \varepsilon_a\ket{\chi_a^\op},
\end{equation}
where $\varepsilon_a$ is the diagonal element $\varepsilon_{aa}$. This is the 
Hartree-Fock equation in its canonical form.
